% s07-sec4.tex —— §4 迹公式的阿代尔化（p.38–46）
\section{迹公式的阿代尔化}
\srcpage{38}

公式 (3.32) 中的函数 $\theta_v(a;s)$ 是 $A_v$ 上紧支的函数。众所周知，它们有界，但不必光滑。我们提供几份参考文献，也就是后续文章中要用的：$\mathbb R$ 与 $\mathbb C$ 上的群，用 Varadarajan 的书 [V]；$p$ 进域上的群，用 Shalika 的文章 [Sh]；正特征局部域上的群，我们没有参考文献。除了这些局部问题之外，乘积
\[
\theta(a;s)=\prod_v\theta_v(a_v;s),\qquad a=\prod_va_v
\]
的行为并不显然，即使取 $s=1$、$f_v$ 为 $G(O_v)$ 的特征函数也是如此。特别地，局部 Fourier 变换 $\hat\theta_v(a;s)$ 不是紧支的，定义整体 Fourier 变换 $\hat\theta(a;s)=\prod_v\hat\theta_v(a;s)$ 的积分之收敛性也未得到。由于我们打算使用 Poisson 公式，$\theta_v$ 与 $\theta$ 的行为需要更细致的研究。本文开始这一研究，对群 $\operatorname{SL}(2)$ 将在下一篇文章中继续。我们希望不久能回到一般问题上来。

通常的阿代尔形式下的 Poisson 公式，是调和分析对自对偶对 $F\subset\mathbb A$ 的一个应用。函数的恶劣行为提示我们使用截断形式，更准确地说，是 Jayce Getz 的一个观察所提示的形式——他在 [G] 中已经遇到过类似的困难。

\srcpage{39}
\begin{lemma}
\setcounter{theorem}{0}
设 $T$ 是整体域 $F$ 上的环面。设 $X$ 是 $T$ 的有理特征群，带有 Galois 群 $\operatorname{Gal}(F^{\mathrm{sep}}/F)$ 的作用。设 $\Lambda=\{\lambda_1,\dots,\lambda_n\}\subset X$ 在 $\operatorname{Gal}(F^{\mathrm{sep}}/F)$ 下不变。最后设 $S$ 是 $F$ 的有限位集，包含所有无穷位，且对每个 $v\in S$ 给定紧子集 $U_v\subset T(F_v)$。则存在有限子集 $S'\supset S$，使得若 $t\in T(F^{\mathrm{sep}})$，$t\in U_v$ 对一切 $v\in S$，且 $|\lambda(t)|_v=1$ 对一切 $v\notin S$ 与一切 $\lambda\in\Lambda$，则
\[
\prod_{j=1}^n\bigl(1-\lambda_j(t)\bigr)=0
\quad\text{或}\quad
|1-\lambda_j(t)|_v=1
\]
对一切 $j$ 与一切 $v\notin S'$ 成立。
\end{lemma}

这个引理虽重要，证明却容易。显然存在正数 $A$，使得对一切 $\lambda\in\Lambda$、$v\in S$ 有 $1/A\le|\lambda(t)|_v\le A$。此外对 $v\notin S$ 有 $|1-\lambda_j|_v\le1$。考虑 $\alpha=\prod_j(1-\lambda_j(t))$。它属于 $F$。对 $v\notin S$，$|\alpha|_v\le1$；对 $v\in S$，$|\alpha|_v\le A^n$。在 $S$ 之外，只有有限个位使 $F_v$ 含有绝对值正且大于或等于 $A^{-n}$ 的元素，也就是使 $q_v\le A^n$ 的位。设 $S'$ 为这些位与 $S$ 的并。集合 $S'$ 满足引理的条件，且只依赖于 $S$、$A$ 与 $n$。显然，当 $S$ 变大或集合 $U_v$（$v\in S$）变大时，$S'$ 会以令人不安的速度变得越来越大。

引理 4.1 允许我们把迹公式重新表述，这对我们很重要。我们从阿代尔 Poisson 公式的一个初等观察开始。设 $S'$ 是 $F$ 的有限位集，包含所有无穷位。固定 $\mathbb A_F$ 的一个整体特征 $\chi$，具有通常的性质：使 $\chi(ba)=1$ 对一切 $a\in F$ 成立的 $b\in\mathbb A_F$ 之集就是 $F$。我们假定 $S'$ 足够大，使得对 $v\notin S'$，$F_v$ 上的自对偶测度给 $O_v$ 的测度为 $1$。于是使 $\chi_v(x)=1$ 的 $x\in F_v$ 之集就是 $O_v$。设 $\mathbb A_F^{S'}$ 是使 $a_v\in O_v$（$v\notin S'$）的 $a\in\mathbb A_F$ 的子集。令
\[
O^{S'}=\prod_{v\notin S'}O_v,\qquad
\mathbb A^{S'}=\prod_{v\in S'}F_v,\qquad
F^{S'}=F\cap\mathbb A_F^{S'}.
\]
我们把 $F^{S'}$ 与它在 $\mathbb A^{S'}$ 中的像等同。

\srcpage{40}
\begin{lemma}
\setcounter{theorem}{1}
若 $S'$ 足够大，则对 $v\notin S'$，$O_v$ 关于 $\chi_v$ 自对偶，且
\begin{equation}\tag{4.3}
\mathbb A_F=\mathbb A_F^{S'}+F.
\end{equation}
此外 $F^{S'}\subset\mathbb A^{S'}$ 关于 $\chi_{S'}$ 在 $\mathbb A^{S'}$ 上的限制自对偶，$\mathbb A^{S'}$ 上的乘积测度给出 $\meas(F^{S'}\backslash\mathbb A^{S'})=1$。
\end{lemma}

显然这个引理对 $F$ 上任何有限维向量空间都有类似的表述。我们将以这一更一般的形式使用它。

所谓 $S'$ 足够大，是指它包含一个我们即将构造的适当的有限子集。为简化记号，我们把这个子集也记作 $S'$。显然等式 (4.3) 是关键，因为它推出
\[
\mathbb A_F=F+O^{S'}+\mathbb A^{S'}.
\]
$\mathbb A^{S'}$ 上在 $F^{S'}$ 上平凡的一个特征，唯一地延拓为 $O^{S'}+\mathbb A^{S'}$ 上在 $O^{S'}+F^{S'}$ 上平凡的特征，进而延拓为 $\mathbb A_F$ 上在 $F$ 上平凡的特征。因此只需对一个 $S'$ 证明引理，然后它对任何更大的集合都成立。

域 $F$ 是某个域 $F^0$ 的有限可分扩张，$F^0$ 或是有理数域，或是有限域上射影直线 $\mathbb P^1$ 的函数域。我们可以取 $S'$ 为 $F^0$ 的某个位集 $S'_0$ 上方的位之集。选取这个集合，使得存在 $F/F^0$ 的一组基 $\{\alpha_1,\dots,\alpha_d\}$，其判别式 $\det(\operatorname{tr}\alpha_i\alpha_j)$ 在 $S_0$ 之外的一切 $v$ 处绝对值为 $1$，且每个 $\alpha_i$ 属于 $F^{S'}$。于是显然
\[
F^{S'}=\Bigl\{\sum_ia_i\alpha_i\ \Big|\ a_i\in F^{0,S'_0}\Bigr\},
\qquad
\mathbb A_F^{S'}=\Bigl\{\sum_ia_i\alpha_i\ \Big|\ a_i\in\mathbb A^{S'_0}_{F^0}\Bigr\}.
\]
因此只需对 $F=F^0$ 证明引理。对 $F=\mathbb Q$，取 $S'=\{\infty\}$ 引理显然成立。对 $\mathbb P^1$ 的函数域，可以取 $S'$ 为单个位，任何一位都行。

迹公式的主部是对椭圆正则共轭类求和。我们已指出，和中的每一项都有一个公因子 $m_G\rho_G\tau(G)$，当 $G$ 半单且单连通时它等于 $1$。为简化公式，以下不再写出这个因子。事实上，本文的主要目的是把迹公式给出的和表示为 Poisson 和，因此此后只考虑半单且单连通的群是合适的。一般情形下要对 $\eta\in h$、对 $B_\eta(F)$ 与对 $Z_\eta(F)$ 求和。造成困难、需要用 Poisson 对偶处理的是对 $B_\eta(F)$ 的求和。参数 $z\in Z_\eta(F)$ 与 $\eta$ 实际上只起指标的作用。因此最好排除它们，而假定 $G=G_{\mathrm{der}}$ 半单且单连通。

\srcpage{41}
由引理 4.1，对函数 $f=\prod f_v$（$f_v$ 光滑紧支），乘积
\[
\prod_v\Orb(\gamma_{\mathrm{st}},f_v)
\]
除有限个稳定正则共轭类外为 $0$。在本文中，一如迹公式的常见情形，奇异类先搁置一旁，等论证的主线更清楚之后再处理。目前还不清楚。因此
\[
\sum_\gamma L(1,\sigma_{T/G})\prod_v\Orb(\gamma_{\mathrm{st}},f_v)
=\lim_{s\downarrow1}\sum_\gamma L(s,\sigma_{T/G})\prod_v\Orb(\gamma_{\mathrm{st}},f_v).
\]
由于对半单正则类有 $|\Delta(\gamma)|=\prod_v|\Delta(\gamma)|_v=1$，这个表达式等于
\begin{equation}\tag{4.4}
\lim_{s\downarrow1}\sum_\gamma\theta\bigl(c(\gamma_{\mathrm{st}});s\bigr).
\end{equation}
和中 $a=c(\gamma_{\mathrm{st}})$ 还只跑遍 $B(F)$ 的一个子集，因为现在只有一个 $\eta$，指标 $\eta$ 已省略。所缺的是非椭圆的 $\gamma_{\mathrm{st}}$。

选取满足引理 4.1 与 4.2 的集合 $S'$。我们还假定在 $S'$ 之外，$G$ 拟分裂且在非分歧扩张上分裂，$f_v$ 是某个超特殊子群的特征函数。第二个引理的条件与公式中使用的函数 $f_v$ 无关；第一个引理的条件则依赖于这些函数的支集。开始分析 Poisson 和时必须考虑到这一点。本文中这些函数与 $S'$ 都固定，但 $S'$ 依赖于 $f_v$ 这一事实无疑会使日后的分析更困难。

我们选取 $S'$，使得若 $v\notin S'$ 且轨道积分 $\Orb(\gamma_{\mathrm{st}},f_v)\ne0$，则 $|\Delta(\gamma)|_v=1$。我们回忆 [O1, O2, O3] 中的一些公式。设 $\kappa_v$ 是 $v$ 的剩余域。首先，对在 $d_v=O_v$ 的位 $v$ 处有好约化的环面 $T$，公式 (3.5) 的右端等于
\begin{equation}\tag{4.5}
q_v^{-\dim T}|T(\kappa_v)|=L(1,\sigma_T)^{-1},
\end{equation}

\srcpage{42}
这个等式就是 [O1] 的公式 (1.2.6)。我们还要验证：在对 $\gamma$ 的假设，即
\begin{itemize}
\item[(i)] $|\lambda(\gamma)|_v=1$ 对 Cartan 子群 $T$ 的每个特征 $\lambda$ 成立，
\item[(ii)] $|1-\alpha(\gamma)|_v=1$ 对每个根 $\alpha$ 成立，
\end{itemize}
之下，有
\begin{equation}\tag{4.6}
\Orb(\gamma_{\mathrm{st}},f_v)=q_v^{-\dim G_{\mathrm{der}}}|G_{\mathrm{der}}(\kappa_v)|.
\end{equation}
这个等式来自假设 (i)、(ii) 对 $\gamma$ 的两个推论：
\begin{itemize}
\item[(a)] 元素 $\gamma$ 本身只含在一个超特殊子群中；
\item[(b)] $\gamma$ 的稳定共轭类与 $G_{\mathrm{der}}$ 的每个超特殊子群 $K$ 都有非空交。
\end{itemize}
这是与 $G(F_v)$ 相联系的仿射大楼的性质的推论。遗憾的是——当然也是由于努力不够——尽管大楼几乎被普遍使用，我们却没有找到恰好给出我们所需要的东西的参考文献。不过我们还是继续下去。设 $f_v$ 是子群 $K$ 的特征函数，$K$ 超特殊。设 $x$ 是 Bruhat–Tits 大楼中定义 $K$ 的点。把 $\gamma$ 换成与它稳定共轭的元素，可以假定它含于 $K$。若 $\gamma'$ 与 $\gamma$ 稳定共轭，则 $\gamma'=g^{-1}\gamma g$，$g\in G(F^{\mathrm{sep}})$。若 $\gamma'$ 固定 $K$ 或 $x$，则由上述三个推论中的第一个（但用于 $F$ 的一个有限扩张）可知 $g$ 固定 $x$，而这时上闭链 $\sigma\to g\sigma^{-1}$ 取值于 $K=G(O_v)$。由 Lang 定理它是平凡的，故 $\gamma'$ 与 $\gamma$ 在 $K$ 中共轭。因此
\[
\Orb(\gamma_{\mathrm{st}},f_v)
=\frac{1}{L_v(1,\sigma_{T/G})}
\frac{\meas^{\mathrm{geom}}K_v}{\meas^{\mathrm{geom}}(K_v\cap T)}
=q_v^{-\dim G_{\mathrm{der}}}|G_{\mathrm{der}}(\kappa_v)|.
\]
我们回忆（[O3]），乘积
\[
\Pi_{S'}=\prod_{v\notin S'}q_v^{-\dim G_{\mathrm{der}}}|G_{\mathrm{der}}(\kappa_v)|
\]
收敛。它显然与 $T$ 无关。我们令
\[
\Pi_{S'}(\gamma_{\mathrm{st}},\mathrm{sd})=L^{S'}(s,\sigma_{T/G})\Pi_{S'}.\yizhu{此处原文第二个变量写作 $sd$，未加解释；对照左侧的 $(\gamma_{\mathrm{st}},s)$，它很可能就是 $s$ 的笔误。译文照原文排印。}{依据：原书 p.42}
\]
考虑迹公式，更确切地说是它的椭圆部分——经过我们的变形，它已成为公式 (4.4)——对给定的函数 $f=\prod f_v$，并选取足够大的集合 $S'$。

\srcpage{43}
由 $S'$ 的选取，对 $v\notin S'$ 有 $|\Delta_v(\gamma)|=1$，于是
\begin{equation}\tag{4.7}
\theta(\gamma_{\mathrm{st}};s)
=L^{S'}(s,\sigma_{T/G})\,\Pi_{S'}\prod_{v\in S'}\theta_v(\gamma_{\mathrm{st}};s).
\end{equation}
我们引入 Getz 截断，只是为了避开一个困难：我们不知道如何验证 $\theta(a;s)$（$s>1$）的 Fourier 变换的行为允许使用 Poisson 对偶。这种可能性不大。等式 (4.7) 又带来另一个困难：因子 $L(s,\sigma_{T/G})$ 使 (4.7) 中出现因子 $L^{S'}(s,\sigma_{T/G})$，而它通过 $T$ 依赖于 $\gamma$。由于对给定的 $f$ 只有有限个 $\gamma$，我们可以取 $S'$ 足够大，使 $L^{S'}(s,\sigma_{T/G})$ 在开区间 $(1,1+\epsilon)$ 上任意接近 $1$，同时不改变 (4.7) 的最后一个因子。第二个因子一致地变化，从而与 $\gamma$ 无关，并当 $S$ 增大时趋于 $1$。由于 (3.25) 中出现的 $\gamma$ 之集不随 $S'$ 增大而改变，我们可以把 (3.17) 中各项之和换成同一个 $\gamma$ 之集上的和
\begin{equation}\tag{4.8}
\Pi_{S'}\prod_{v\in S'}\theta_v(\gamma_{\mathrm{st}};1)
=\Pi_{S'}\lim_{s\downarrow1}\prod_{v\in S'}\theta_v(\gamma_{\mathrm{st}};s).
\end{equation}
考虑到 $\Pi_{S'}$ 趋于 $1$，出于同样的理由我们也可以把它去掉。\footnote{在此我们对"阿代尔化"一词的用法作一简短说明。数域 $F$ 上的向量空间 $V$ 及其张量积 $\mathbb A_F\times V$，在解析数论中用来借助 Poisson 公式研究 $\sum_{v\in V}\phi(v)$（$\phi=\prod_v\phi_v$）及其渐近行为。这个公式是仅有的几种可用工具之一。我们想把它用于 $\theta(a;s)$。这是不允许的，因为函数 $\theta_v$ 不光滑。但它们的性状看来足够好，使得对它们的使用不是在局部域的受限无穷乘积中进行，而是在截断乘积中进行。遗憾的是，我们无法在几页纸内解释一切。老实说，我们自己也几乎没有开始研究这些问题（[L5]）。}

不过我们回避了一个小困难：定义 $L^{S'}(s,\sigma_{T/G})$ 的无穷乘积在区间上不一致收敛。一致收敛的是外层的乘积
\[
\prod_{n=1}^\infty\Bigl(\prod_{n<q_v\le n+1}L_v(s,\sigma_{T/G})\Bigr).
\]
这是所谓 Chebotarev 定理的一种形式，[IK] 中有一种形式。我们没有再找其他参考文献。这样的文献很多。\srcpage{44}
因此我们始终假定 $S'$ 是形如 $\{v\mid q_v\le n\}$ 的集合，这并不妨碍它足够大。

为研究 (4.8) 中各项之和——即迹公式中出现、但删去了 $\Pi_{S'}$ 的和
\begin{equation}\tag{4.9}
\sum_\gamma\prod_{v\in S'}\theta_v(\gamma_{\mathrm{st}};s),
\end{equation}
我们可以对 $F^{S'}\subset\mathbb A^{S'}$ 使用 Poisson 对偶。事实上，由引理 4.1 与 $S'$ 的选取，若对一切 $v\in S'$ 有 $\theta_v(\gamma_{\mathrm{st}},1)\ne0$，则 $c(\gamma_{\mathrm{st}})$ 在 $S'$ 之外为整数。于是和 (4.9) 可写成
\begin{equation}\tag{4.10}
\sum_{b\in B^{S'}}\prod_{v\in S'}\theta_v(b;s),
\end{equation}
其中 $B^{S'}$ 定义为 $\operatorname{Spec}F^{S'}[b_1,\dots,b_r]$。

我们关于可以使用 Poisson 对偶的说法是乐观的，而且还不完全正确。首先函数 $\theta_v(b;s)$ 不光滑，它们在子簇 $D(b)=0$ 上奇异。不过它们很可能足够光滑，也就是说奇异性足够温和，使得对偶 Poisson 和
\begin{equation}\tag{4.11}
\sum_{b\in B^{S'}}\prod_{v\in S'}\hat\theta_v(b;s)
\end{equation}
绝对收敛，并等于函数 $\prod_{S'}\theta_v(b;s)$ 的 Poisson 和。我们中的一人希望在下一篇文章中对 $\operatorname{SL}(2)$ 验证这一点（[L5]）。没有理由认为即使对任意群结果也会如此困难。所需理论的基础很可能在 [Sh, V] 中。以后有更精确的记号会方便些：
\begin{equation}\tag{4.12}
\theta_{S'}(b;s)=\prod_{v\in S'}\theta_v(b;s),
\qquad
\hat\theta_{S'}(b;s)=\prod_{v\in S'}\hat\theta_v(b;s).
\end{equation}
第二个困难是，我们从一开始就排除了不定义椭圆环面的 $b$。我们可以把它们加进 (4.10)。但随后又必须减去，因为它们带来收敛问题。特别地，对它们来说 $L^{S'}(s,\sigma_{T/G})$ 当 $S'$ 越来越大、$s\downarrow1$ 时趋于无穷。不过可以设想这些不同的无穷会相互抵消。例如对群 $\operatorname{SL}(2)$，在考察形如 (1.4) 的和的性状时，自然会减去二面体群表示的转移所带来的贡献，甚至加上分裂环面，\srcpageinl{45}这使加减过程更自然。尽管如此，仍有一个头疼的问题：在这样得到的公式中，两个无穷究竟如何相互抵消。这个问题将在 [L5] 中重新讨论。

为此必须引入并研究稳定转移 $f=f^G\mapsto f^H$ 与转移因子的概念。最后要研究的不是单个函数 $\theta$ 或 $\hat\theta$ 的 Poisson 和，而是差 $\theta-\phi$ 与 $\hat\theta-\hat\phi$ 的 Poisson 和，其中 $\phi$ 由转移 $f^H$ 定义。即使没有 $\phi$，我们也可以考察当 $s\to1$ 时仍然有限的项 $\hat\theta(0;s)$，即使加上与非椭圆环面相联系的项也是如此。我们将看到（但不在本文中）$\hat\phi(0;s)$ 的贡献为 $0$。

在开始之前，回顾一下这一系列变形。我们引入参数 $s>1$，并须让它趋于 $1$。然后引入足够大的集合 $S'$，过渡到截断和。对这个截断和使用 Poisson 对偶。但这时对偶和中的项
\[
\hat\theta_{S'}(0)=\lim_{s\to1}\hat\theta_{S'}(0;s)
\]
依赖于 $S'$。为得到与 $S'$ 无关的值，必须取二重极限
\begin{equation}\tag{4.13}
\hat\theta(0)=\lim_{S'\to\infty,\ s\downarrow1}\hat\theta_{S'}(0;s),
\end{equation}
其中 $S'$ 按所规定的方式变得越来越大。本文最后一节讨论的正是这个极限值。不仅对主项，而且对剩下的和 $\sum_{b\ne0}\hat\theta_{S'}(b;s)$ 也必须取这个二重极限。对这个和，要处理的不是它本身，而是差
\[
\sum_{b\ne0}\hat\theta_{S'}(b;s)-\sum_{b\ne0}\hat\phi_{S'}(b;s).
\]
对项 $\hat\phi_{S'}(0;s)$ 的研究将只在下一篇文章中进行。\footnote{一位审稿人对这个二重极限感到困惑。他有理由如此。事实上，极限的次序无关紧要。这些极限，特别是 $s\downarrow1$ 的极限，将在 [L5] 中更仔细地讨论。}
